\chapter{MỞ ĐẦU}

\section{Giới thiệu bài toán Object Detection}

Phát hiện vật thể (Object Detection) là một trong những bài toán trung tâm của thị giác máy tính (Computer Vision). Mục tiêu là xác định \textbf{vị trí} (bounding box) và \textbf{loại} (class) của tất cả vật thể trong ảnh hoặc video. Object Detection được ứng dụng rộng rãi trong xe tự hành, giám sát an ninh, y tế, công nghiệp, nông nghiệp và nhiều lĩnh vực khác.

Trước khi YOLO ra đời (2015), các phương pháp detection chủ đạo là \textbf{two-stage detectors}:
\begin{itemize}
    \item \textbf{R-CNN} (2014): Dùng Selective Search tạo ~2000 region proposals, mỗi region qua CNN trích features, rồi SVM phân loại. Cực chậm: >40 giây/ảnh.
    \item \textbf{Fast R-CNN} (2015): Cải tiến bằng cách chia sẻ CNN features, nhưng vẫn cần Selective Search.
    \item \textbf{Faster R-CNN} (2015): Thay Selective Search bằng Region Proposal Network (RPN), nhưng vẫn là pipeline 2 giai đoạn.
\end{itemize}

Điểm chung của các phương pháp trên: pipeline \textbf{rời rạc}, mỗi thành phần tối ưu riêng, khó đạt real-time.

\section{YOLO — Bước ngoặt}

YOLO (You Only Look Once) ra đời năm 2015, thay đổi hoàn toàn cách tiếp cận: coi object detection là \textbf{bài toán hồi quy duy nhất}, từ pixel ảnh trực tiếp đến bounding boxes và class probabilities. Chỉ cần \textbf{một lần forward pass} qua mạng neural → real-time (45 FPS).

Từ 2015 đến 2026, YOLO đã trải qua \textbf{14 phiên bản chính} (v1--v12, YOLOE, YOLO26) và nhiều biến thể (YOLOX, PP-YOLOE, YOLO-NAS, YOLO-World), với mAP tăng từ 63.4\% (VOC) lên ~57.2\% (COCO), tốc độ từ 45 FPS lên hàng trăm FPS.

\section{Mục tiêu và phạm vi báo cáo}

\textbf{Mục tiêu:} Khảo sát toàn diện sự phát triển của YOLO qua các đời, phân tích kiến trúc, phương pháp, kết quả, ưu/nhược điểm, và xu hướng tương lai.

\textbf{Phạm vi:}
\begin{itemize}
    \item Dòng chính: YOLOv1 → YOLOv12, YOLOE, YOLO26
    \item Biến thể: YOLOX, PP-YOLOE, YOLO-NAS, YOLO-World
    \item So sánh đa chiều và khuyến nghị lựa chọn
\end{itemize}
